\chapter*{ABSTRACT}
\addcontentsline{toc}{chapter}{Abstract}
\begingroup
\small
\setlength{\parskip}{1.5mm}

Every year a large cohort of Bangladeshi students prepares for the engineering
university admission tests (BUET, KUET, RUET, CUET) by working through
undifferentiated question banks. A student who repeatedly fails questions on
integration by parts is rarely told that the actual gap is algebraic
manipulation two levels below; coaching centres cannot diagnose at that
granularity at scale, and existing digital platforms largely serve the same
fixed content to everyone. The engineering problem is therefore one of
inference under uncertainty: from a sparse, noisy sequence of multiple-choice
responses, locate a learner's knowledge frontier precisely enough to act on it.

This report presents ATLAS, an adaptive tutoring system that models a learner's
latent mastery with Bayesian Knowledge Tracing (BKT) over a directed acyclic
graph (DAG) of prerequisite relations between granular skills. Three extensions
to textbook BKT carry the design. First, the guess and slip parameters are
conditioned on the Bloom's Taxonomy level of the question, so that a correct
answer on a recall item and a correct answer on an analysis item update belief
by different amounts. Second, a correct response propagates upward through the
prerequisite graph, raising every ancestor to at least the descendant's
posterior, since demonstrating a dependent skill is evidence for the skills it
depends on. Third, a skill's learning-rate parameter is gated conjunctively:
it rises only when \emph{all} of its direct prerequisites are mastered,
preventing a learner from appearing to advance past an unmastered foundation.
An adaptive diagnostic selects each next question by a balanced-volume
heuristic over untested ancestors and descendants, at one Bloom level above the
learner's current estimated band.

The system was built as a FastAPI backend, a React frontend rendering Bangla
text with embedded \LaTeX{} mathematics, and a PostgreSQL (Supabase) store. A
retrieval-augmented generation pipeline over six HSC textbooks (4{,}474 parsed
question records) produces the question bank, with a two-model
generate-then-verify design in which a second model blind-re-solves each
generated item before it is accepted. The knowledge structure itself was the
dominant engineering effort: an initial 430-skill ontology derived from a
sample of the corpus proved too sparse to support the gating policy, and was
rebuilt from the full corpus into 1{,}688 skills and 1{,}743 prerequisite edges
covering all 183 syllabus topics, validated as an acyclic graph of depth 9 with
no isolated skills.

Verification is structural and functional rather than empirical: an 18-check
end-to-end harness exercises the diagnostic, mastery propagation, the
ancestor-ordering invariant and prerequisite spillover against an in-memory
database stand-in, and passes completely; the ontology passes an automated
structural validator with zero errors. The system has not yet been evaluated
with human learners, and establishing that its mastery estimates predict real
performance remains the principal item of future work.

\vspace{12pt}
\noindent\textbf{Keywords:} Bayesian Knowledge Tracing; adaptive learning;
prerequisite graph; Bloom's Taxonomy; intelligent tutoring systems;
retrieval-augmented question generation
\endgroup
\clearpage
